\documentclass[10pt,a4paper]{amsart}
\usepackage[margin=19mm]{geometry}
\usepackage{amsmath,amssymb,mathtools}
\usepackage[scr=rsfs,cal=euler]{mathalfa}
\usepackage{xfrac}
\usepackage[unicode,hidelinks,bookmarks=false]{hyperref}
\newcommand{\ie}{i.e.\ }
\newcommand{\cf}{cf.\ }
\newcommand{\resp}{respectively\ }
\newcommand{\Subsection}{\subsection}
\providecommand{\nobreakdash}{\nobreak\hbox{-}}
\setlength{\emergencystretch}{2em}
\multlinegap0pt
\usepackage{bm}
\usepackage{bbm}
\usepackage{dsfont}
\usepackage{xspace}
\usepackage{cancel}
\usepackage{mathtools}
\usepackage{amsthm}
\makeatletter
\@ifundefined{theorem}{\newtheorem{theorem}{Theorem}}{}
\@ifundefined{lemma}{\newtheorem{lemma}[theorem]{Lemma}}{}
\makeatother
\DeclareMathOperator{\Vol}{Vol}
\DeclareMathOperator{\dvol}{dV}
\newcommand{\bB}{\mathbb{B}}
\newcommand{\bR}{\mathbb{R}}
\newcommand{\bS}{\mathbb{S}}
\newcommand{\mM}{\mathcal{M}}
\title[Energy reduction for Fourth order Willmore energy]{Energy reduction for Fourth order Willmore energy}
\author{Nan Wu, Zetian Yan}
\date{Source: arXiv:2504.08105v3 (CC BY 4.0)}
\begin{document}
\maketitle

\noindent\emph{Source and licence.} Energy reduction for Fourth order Willmore energy. Version arXiv:2504.08105v3 (CC BY 4.0), distributed under the Creative Commons Attribution 4.0 International licence, as recorded in the arXiv licence metadata for this exact version and shown on the abstract page https://arxiv.org/abs/2504.08105v3. Full rights notice: licenses/arxiv-2504.08105-CC-BY-4.0.txt.\par\smallskip

\noindent\textit{Editorial scope.} Counted unit: the Lemma whose source label is \texttt{energy-f} in the article (source0.tex lines 2126--2142, source label \texttt{energy-f}), delivered together with the article's own complete original proof of it (source0.tex lines 2143--2245, one \texttt{proof} environment of 103 lines). No mathematics is written, repaired or re-derived: statement and proof are the article's own text line for line. Required context retained uncounted: f (lines 1382--1407). Every definition, display and label of those lines is retained unchanged. Omitted from this package and not redistributed: 1--1381, 1408--2125, 2246--2795. Nothing inside a retained excerpt is omitted or altered: every retained body is delivered exactly as the article's lines are written, so no omission receipt of any kind accompanies this package. Citations: the retained excerpts carry no citation command at all, so no bibliography entry of the article is transcribed and no reference is redistributed. Reference adaptation: none was needed and none was made; every reference inside the delivered unit resolves inside it, so no unresolved reference or citation is produced. Printed slips kept literal: no printed word, symbol, hypothesis or number of the counted statement or of its proof was altered. Layout and class: the article's own class is \texttt{amsart} with packages including amsmath, amssymb, amsthm, amsrefs, hyperref, cleveref, mathrsfs, mathtools, slashed, tikz-cd, enumitem; the wrapper is the lane's pinned \texttt{amsart} editorial wrapper at 10pt on A4, which restates the theorem-like environments the retained text needs and adds the article's own macro definitions that the retained text needs (\texttt{\textbackslash Vol}, \texttt{\textbackslash bB}, \texttt{\textbackslash bR}, \texttt{\textbackslash bS}, \texttt{\textbackslash dvol}, \texttt{\textbackslash mM}), with extra wrapper packages (bm, bbm, dsfont, xspace, cancel, mathtools). The delivered numbering is the unit's own. Assets: no retained span contains a figure environment, a graphics-inclusion command, a PDF-page inclusion, a table or a TikZ picture; no image file is copied into this package, the package declares no asset dependency and no image file is redistributed with it. Licence: version arXiv:2504.08105v3 of this article is distributed under the Creative Commons Attribution 4.0 International licence, as recorded in the arXiv licence metadata for this exact version and shown on the abstract page https://arxiv.org/abs/2504.08105v3. A full rights notice is delivered at licenses/arxiv-2504.08105-CC-BY-4.0.txt. Characters: every retained excerpt is pure ASCII, so the delivered engine, XeLaTeX with its default Latin Modern text font, reports no missing character.
\begin{lemma}\label{f}
    If $f(r)$ is a radial function satisfying $\Delta_0^3 (f(r)X_i)=0$, for some $X_i\in \mathcal{H}^{\bS}_{1}$, $f(r)$ must be in the form of
    \begin{equation}
        f(r)=\sum_{i=1}^{6}B_i F_i(r), \quad B_i\in \bR^{n-4},
    \end{equation}
where
\begin{equation*}
    \begin{split}
        &F_1(r)=r^{-3}, \quad F_2(r)=r^{-1}, \quad F_3(r)=r,\\
        &F_4(r)=r\ln r, \quad F_5(r)=r^3, \quad F_6(r)=r^5.
    \end{split}
\end{equation*}
Moreover, the matrix $\mM^{f}$ encodes values of $\{F_i(r)\}_{i=1}^{6}$ on $r=\gamma$ and $r=1$:
\begin{equation}\label{Mf}
\mM^{f}:=\begin{bmatrix}
\gamma^{-3} & \gamma^{-1} &  \gamma & \gamma \ln \gamma & \gamma^3& \gamma^5\\
-3\gamma^{-4} & -\gamma^{-2} & 1 & 1+\ln \gamma & 3\gamma^2& 5\gamma^4\\
12\gamma^{-5} & 2\gamma^{-3} & 0 & \gamma^{-1} & 6 \gamma& 20\gamma^3\\
1 & 1 & 1 & 0 & 1& 1\\
-3 & -1 & 1 & 1 & 3& 5\\
12 & 2 & 0 & 1 & 6& 20
\end{bmatrix}.
\end{equation}


\end{lemma}

\begin{lemma}\label{energy-f}
    Let $F(r)$ be the vector defined in Lemma \ref{f}. Then for each $X_i\in \mathcal{H}^{\bS}_{1}$, we have
\begin{equation}
\begin{split}
    &\int_{\bB_{\sigma}(0)\backslash \bB_{\tau}(0)}\langle \nabla \Delta_0 (F(r)\cdot X_i), \nabla \Delta_0(F(r)\cdot X_i)\rangle \dvol\\
    =&16 \Vol(\bS^3) \begin{bmatrix}
0 & 0 & 0 & 0 & 0& 0\\
0 & -3r^{-4} & 0 & 3r^{-2} & 0& 24r^2\\
0 & 0 & 0& 0& 0 & 0\\
0 & 3r^{-2} & 0 & {4\ln r} & 3r^2& 0\\
0 & 0 & 0 & 3r^2 & 9r^4& 24r^6\\
0 & 24r^2 & 0 & 0 & 24r^6& 96r^8
\end{bmatrix}\Bigg|_{r=\tau}^{r=\sigma}.
\end{split}
\end{equation}
    
\end{lemma}

\begin{proof}
    From Lemma \ref{f}, we already know that 
    the Laplace operator $\Delta_0$ acts on the six-dimensional space $F(r)\cdot X_i$ as the matrix:
\begin{equation}
    \Delta_0(F(r)\cdot X_i)=F(r)\cdot X_i\cdot 
    \begin{bmatrix}
0 & -4 & 0 & 0 & 0& 0\\
0 & 0 & 0 & 4 & 0& 0\\
0 & 0 & 0 & 0 & 12& 0\\
0 & 0 & 0 & 0 & 0& 0\\
0 & 0 & 0 & 0 & 0& 32\\
0 & 0 & 0 & 0 & 0& 0
\end{bmatrix},
\end{equation}
and
\begin{equation*}
        \partial_r \left(F(r)\cdot X_i\right)=r^{-1}F(r)\cdot X_i\cdot\begin{bmatrix}
-3 & 0 & 0 & 0 & 0& 0\\
0 & -1 & 0 & 0 & 0& 0\\
0 & 0 & 1 & 1 & 0& 0\\
0 & 0 & 0 & 1 & 0& 0\\
0 & 0 & 0 & 0 & 3& 0\\
0 & 0 & 0 & 0 & 0& 5
        \end{bmatrix}.
    \end{equation*}

The bi-Laplace operator $\Delta^2_0$ acts on the six-dimensional space $F(r)$ as the matrix:
\begin{equation}
    \Delta^2_0(F(r)\cdot X_i)=F(r)\cdot X_i\cdot 
    \begin{bmatrix}
0 & 0 & 0 & -16 & 0& 0\\
0 & 0 & 0 & 0 & 0& 0\\
0 & 0 & 0 & 0 & 0& 384\\
0 & 0 & 0 & 0 & 0& 0\\
0 & 0 & 0 & 0 & 0& 0\\
0 & 0 & 0 & 0 & 0& 0
\end{bmatrix}.
\end{equation}

Therefore, we have
\begin{equation}
\begin{split}
    &-\int_{\bB_{\sigma}(0)\backslash \bB_{\tau}(0)}\langle \Delta_0 (F(r)\cdot X_i), \Delta^2_0(F(r)\cdot X_i)\rangle \dvol\\
=&-64\int_{\bS^3} X^2_i\dvol_{g_c}\int_{\tau}^{\sigma}\begin{bmatrix}
0 & 0 & 0 & 0 & 0& 0\\
0 & 0 & 0 & r^{-6} & 0& -24r^{-2}\\
0 & 0 & 0 & 0 & 0& 0\\
0 & 0 & 0 & -r^{-4} & 0& 24\\
0 & 0 & 0 & -3r^{-2} & 0& 72r^2\\
0 & 0 & 0 & -8 & 0 & 192r^4
\end{bmatrix}r^3dr\\
=&-64\Vol(\bS^3) \begin{bmatrix}
0 & 0 & 0 & 0 & 0& 0\\
0 & 0 & 0 & -\frac{r^{-2}}{2} & 0& -12r^2\\
0 & 0 & 0 & 0 & 0& 0\\
0 & 0 & 0 & -\ln r & 0& 6r^4\\
0 & 0 & 0 & -\frac{3}{2}r^2 & 0& 12r^6\\
0 & 0 & 0 & -2r^4 & 0 & 24r^8
\end{bmatrix}\Bigg|_{r=\tau}^{r=\sigma}.
    \end{split}
\end{equation}
Besides,
\begin{equation}
    \partial_r \circ \Delta_0(F(r)\cdot X_i)=r^{-1} F(r)\cdot X_i\cdot 
    \begin{bmatrix}
0 & 12 & 0 & 0 & 0& 0\\
0 & 0 & 0 & -4 & 0& 0\\
0 & 0 & 0 & 0 & 12& 0\\
0 & 0 & 0 & 0 & 0& 0\\
0 & 0 & 0 & 0 & 0& 96\\
0 & 0 & 0 & 0 & 0& 0
\end{bmatrix}.
\end{equation}
We have
\begin{equation}
\begin{split}
    &\int_{\partial \bB_{\sigma}(0)\backslash \partial \bB_{\tau}(0)} \langle \Delta_0(F(r)\cdot X_i), \partial_r \circ \Delta_0(F(r)\cdot X_i) \rangle dA\\
    =&16 \Vol(\bS^3) \begin{bmatrix}
0 & 0 & 0 & 0 & 0& 0\\
0 & -3r^{-4} & 0 & r^{-2} & -3& -24r^2\\
0 & 0 & 0& 0& 0 & 0\\
0 & 3r^{-2} & 0 & -1 & 3r^2& 24r^4\\
0 & 9 & 0 & -3r^2 & 9r^4& 72r^6\\
0 & 24r^2 & 0 & -8r^4 & 24r^6& 192r^8
\end{bmatrix}\Big|_{r=\tau}^{r=\sigma}.
    \end{split}
\end{equation}
Combining them together yields
\begin{equation}
\begin{split}
    &\int_{\bB_{\sigma}(0)\backslash \bB_{\tau}(0)}\langle \nabla \Delta_0 (F(r)\cdot X_i), \nabla \Delta_0(F(r)\cdot X_i)\rangle \dvol\\
    =&16 \Vol(\bS^3) \begin{bmatrix}
0 & 0 & 0 & 0 & 0& 0\\
0 & -3r^{-4} & 0 & 3r^{-2} & 0& 24r^2\\
0 & 0 & 0& 0& 0 & 0\\
0 & 3r^{-2} & 0 & {4\ln r} & 3r^2& 0\\
0 & 0 & 0 & 3r^2 & 9r^4& 24r^6\\
0 & 24r^2 & 0 & 0 & 24r^6& 96r^8
\end{bmatrix}\Big|_{r=\tau}^{r=\sigma},
\end{split}
\end{equation}
Here we have replaced constant functions $-3$ and $9$ by zeros which will not affect our results.
\end{proof}

\end{document}
